\chapter{Insect Bites and Stings}
\index{Insect Bites and Stings}

\section*{Definition and Overview}
Insect bites and stings are common in many countries, where many insects, spiders, and other arthropods are active year-round. Bites from mosquitoes, midges, and fleas cause itchy lumps; stings from bees, wasps, and ants inject venom and can cause pain, swelling, and, in allergic people, severe reactions. the affected region also has some of the world's most dangerous spiders, including the funnel-web and redback, although deaths are now very rare because effective antivenoms exist. Most bites and stings are managed at home with simple first aid, but anaphylaxis, snake-like spider bites, and multiple stings require urgent medical care. Knowing the correct first aid for local creatures can save lives.

\section*{Epidemiology}
Mosquito bites are almost universal in many countries, and mosquito-borne diseases such as Ross River virus and Murray Valley encephalitis occur in some regions. Bee and wasp stings send many thousands of people to hospital each year, and anaphylaxis from stings causes around a dozen deaths annually. Redback spider bites are common and cause painful symptoms, while funnel-web spider bites are rare but dangerous, occurring mainly in New South Wales and southern Queensland. Tick bites are increasing, particularly along the east coast, and can transmit infections or cause allergy to meat. Children, outdoor workers, and people in rural and coastal areas are most exposed.

\section*{Aetiology and Pathophysiology}
An insect bite is caused by the insect puncturing the skin to feed, injecting saliva that contains anticoagulants and proteins, to which the body mounts an immune and inflammatory response producing itching, redness, and swelling. Stings deliver venom through a stinger; the venom's toxins cause pain and local tissue damage, and in susceptible people trigger allergic reactions ranging from large local swelling to anaphylaxis, a severe, life-threatening reaction involving the airway, breathing, and circulation. Spider venoms vary: redback venom affects nerve endings causing severe pain and autonomic symptoms, while funnel-web venom is neurotoxic and rapidly dangerous. Tick paralysis results from a toxin in tick saliva.

\section*{Clinical Features}
\begin{itemize}
    \item Mosquito, midge, and flea bites: itchy red lumps that settle in days
    \item Bee and wasp stings: immediate sharp pain, redness, and swelling at the site
    \item Large local reactions: extensive swelling lasting days
    \item Anaphylaxis: difficulty breathing, swelling of the face and throat, wheezing, collapse, or hives over the body
    \item Redback spider bite: severe, spreading pain, sweating, muscle weakness, and nausea
    \item Funnel-web spider bite: severe pain, agitation, salivation, and rapid deterioration
    \item Tick bites: local irritation, and occasionally paralysis or allergy to meat
    \item Infected bites: increasing pain, redness, pus, and fever
\end{itemize}

\section*{Differential Diagnosis}
Skin reactions after an insect bite may be confused with other conditions.
\begin{itemize}
    \item \textbf{Cellulitis and skin infection:} worsening redness, warmth, and pain
    \item \textbf{Allergic reactions:} hives and angioedema from foods, medicines, or stings
    \item \textbf{Contact dermatitis:} reactions to plants or chemicals
    \item \textbf{Skin infestations:} scabies and lice causing itching
    \item \textbf{Viral exanthems:} rashes in children
    \item \textbf{Localised pain syndromes:} muscle strain or injury mistaken for a bite
    \item \textbf{Systemic infection from bites:} mosquito-borne viruses and rickettsial disease
\end{itemize}

\section*{When to Seek Medical Care}
Most bites and stings are managed at home, but medical care is needed for severe pain, spreading infection, allergy symptoms, bites from dangerous spiders, and tick bites with paralysis or persistent illness. People with a known severe allergy to stings should carry their adrenaline autoinjector and action plan, and should seek medical care after any significant reaction.

\textbf{Contact your local emergency services for:}
\begin{itemize}
    \item Signs of anaphylaxis: difficulty breathing, throat swelling, wheezing, hoarseness, or collapse
    \item A funnel-web spider bite, or a suspected bite in an area where they live
    \item Collapse, confusion, or severe illness after any bite or sting
    \item Multiple stings, especially in a child
    \item A bite or sting to the mouth or throat
\end{itemize}

\section*{Diagnosis and Investigations}
\begin{itemize}
    \item \textbf{History and examination:} the creature if seen, the site, and the pattern of symptoms
    \item \textbf{Identification of the creature:} photographs or the spider itself help decide management
    \item \textbf{Assessment of allergy:} the severity of the reaction
    \item \textbf{Blood tests:} for infections such as Ross River virus or rickettsial disease where relevant
    \item \textbf{Allergy testing:} skin or blood testing after significant sting reactions
    \item \textbf{Tick removal and testing:} careful removal with fine forceps and consideration of tick-borne illness
\end{itemize}

\section*{Management}
\begin{itemize}
    \item \textbf{Bee stings:} scrape the stinger out, apply a cold pack, and take simple analgesics and antihistamines
    \item \textbf{Wasp and ant stings:} cold pack, elevation, and symptomatic treatment
    \item \textbf{Mosquito and other bites:} cold packs, antihistamines, and soothing creams
    \item \textbf{Anaphylaxis:} adrenaline given promptly by autoinjector, followed by urgent hospital care
    \item \textbf{Redback spider bite:} cold pack and hospital assessment; antivenom if symptoms are significant
    \item \textbf{Funnel-web spider bite:} pressure immobilisation bandage, keep the person still, and urgent hospital care with antivenom
    \item \textbf{Ticks:} careful removal without squeezing, and medical review for symptoms
    \item \textbf{Infected bites:} cleaning and antibiotics for secondary infection
\end{itemize}

\section*{Complications}
\begin{itemize}
    \item Anaphylaxis and death from severe allergic reactions
    \item Systemic effects of venom, including neurotoxicity and muscle damage
    \item Secondary bacterial infection of bites and stings
    \item Mosquito-borne diseases: Ross River virus, dengue, and others
    \item Tick paralysis and meat allergy after tick bites
    \item Chronic itching, scarring, and skin discolouration
    \item Anxiety and fear of insects affecting outdoor life
\end{itemize}

\section*{Prognosis and Outcomes}
Almost all bites and stings resolve without specific treatment. Anaphylaxis is life-threatening but responds dramatically to early adrenaline, and deaths from stings are preventable with prompt treatment. Deaths from spider bites in many countries are extremely rare because antivenoms are highly effective when given early. Most venom effects settle within hours to days, though redback pain can persist. Education about correct first aid and carrying adrenaline for allergic people are the most important measures.

\section*{Prevention and Self-Care}
\begin{itemize}
    \item Wear protective clothing and closed shoes in bushland and high-risk areas
    \item Use insect repellent containing DEET or picaridin
    \item Keep food covered outdoors to reduce wasp and bee attraction
    \item Do not disturb nests or hives
    \item Shake out clothing, boots, and bedding before use
    \item Check for ticks after walking in tick areas and remove them properly
    \item If severely allergic, carry an adrenaline autoinjector and an action plan
    \item Learn the correct first aid for local spiders and stings
    \item Keep the yard free of rubbish and standing water
    \item Seek prompt care for severe reactions, and for bites that become infected
\end{itemize}

\section*{Sources and Further Reading}
The following reputable sources provide further information for healthcare students and patients seeking reliable, current advice about insect bites and stings.
\begin{itemize}
    \item Australian Venom Research Unit. \textit{First aid and bite information.} https://www.avru.org/
    \item Healthdirect Australia. \textit{Insect bites and stings.} https://www.healthdirect.gov.au/
    \item ASCIA. \textit{Allergy and anaphylaxis information.} https://www.allergy.org.au/
    \item Queensland Health. \textit{Spider and insect bites.} https://www.health.qld.gov.au/
    \item St John Ambulance Australia. \textit{First aid for bites and stings.} https://stjohn.org.au/
    \item NHS. \textit{Insect bites and stings.} https://www.nhs.uk/conditions/insect-bites-and-stings/
\end{itemize}
